% ==============================================================================
% config/fonts.tex
% ==============================================================================
% پیکربندی فونت‌ها و زی‌پرشین. 
% بسیار مهم: پکیج xepersian باید همواره آخرین پکیجی باشد که در سند لود می‌شود!

\usepackage{xepersian}

% اطمینان از اعمال شدن صفحه عنوان سفارشی
\AtBeginDocument{
    \let\maketitle\makecover
}

% ------------------------------------------------------------------------------
% تنظیمات فونت‌ها
% ------------------------------------------------------------------------------
% فونت اصلی متن فارسی. 
% نکته: می‌توانید نام فونت را به Vazirmatn، B Nazanin، یا XB Niloofar تغییر دهید.
\settextfont[Scale=1.1, AutoFakeBold, AutoFakeSlant]{Vazirmatn}

% فونت بخش‌های لاتین سند
\setlatintextfont[Scale=1.0]{Times New Roman}

% تنظیم فونت اعداد به صورت فارسی در متن
\setdigitfont[Scale=1.1, AutoFakeBold, AutoFakeSlant]{Vazirmatn}

% تنظیم فونت اعداد در محیط‌های ریاضی به Yas
\setmathdigitfont[Scale=1.1]{Yas}

% فونت مونو اسپیس برای نمایش کدهای برنامه‌نویسی و مسیرها
\setmonofont[Scale=0.9]{Consolas}

% ------------------------------------------------------------------------------
% تنظیمات فاصله‌ها
% ------------------------------------------------------------------------------
\linespread{1.3} % ضریب فاصله بین خطوط (تقریباً معادل 1.5 در MS Word)
\setlength{\parskip}{0.5\baselineskip} % فاصله کمی بیشتر بین پاراگراف‌ها برای خوانایی بهتر
